%
% @file doc.tex
% @author David Kvaček (xkvace00@stud.fit.vutbr.cz)
% @brief Source code of the project report from the IZV course.
% @date 2024-01-03
%

\documentclass[a4paper, 10pt]{article}
\usepackage[left=2cm, text={18cm, 26cm}, top=2cm]{geometry}
\usepackage[czech]{babel}
\usepackage[utf8]{inputenc}
\usepackage[IL2]{fontenc}
\usepackage{graphicx}
\usepackage{hyperref}
\usepackage{caption}
\pagestyle{empty}

\begin{document}

\begin{center}
    \section*{Vlastní analýza dopravních nehod}
    Zpráva třetí části projektu v~rámci kursu IZV \\
    Autor: David Kvaček (\href{mailto:xkvace00@stud.fit.vutbr.cz}{xkvace00@stud.fit.vutbr.cz})
\end{center}

Cílem této analýzy je zkoumání vztahů mezi hlavními příčinami dopravních nehod, průměrnou hmotnou škodou a~vlivem alkoholu či drog na viníky nehod.
Analýza je provedena na základě poskytnutých dat o~dopravních nehodách v~letech 2023 až 2024.

Pro analýzu jsou data nejprve připravena, což zahrnuje jejich načtení a~předzpracování z~poskytnutých souborů, následované filtrací neplatných nebo chybějících hodnot v~příslušných sloupcích.
Hlavní příčiny nehod jsou rozděleny do šesti kategorií: nezaviněná řidičem, nepřiměřená rychlost jízdy, nesprávné předjíždění, nedání přednosti v~jízdě, nesprávný způsob jízdy a~technická závada vozidla.
Pro výpočet vlivu alkoholu a~drog jsou řidiči klasifikováni podle přítomnosti těchto látek v~krvi do čtyř skupin: bez vlivu alkoholu a~drog, pod vlivem alkoholu, pod vlivem drog a~pod vlivem alkoholu a~drog.

\begin{figure}[ht]
    \centering
    \includegraphics[width=0.86\linewidth]{fig.png}
\end{figure}

\begin{table}[ht]
    \caption*{\textbf{Počet nehod podle hlavní příčiny a~vlivu alkoholu a~drog}}
    \resizebox{\textwidth}{!}{
        \begin{tabular}{|l|r|r|r|r|}
        \hline
        \textbf{hlavní příčina nehody} &
          \multicolumn{1}{l|}{\textbf{bez vlivu alkoholu a~drog}} &
          \multicolumn{1}{l|}{\textbf{pod vlivem alkoholu}} &
          \multicolumn{1}{l|}{\textbf{pod vlivem drog}} &
          \multicolumn{1}{l|}{\textbf{pod vlivem alkoholu a~drog}} \\ \hline
            nezaviněná řidičem         & \texttt{807}   & \texttt{103}  & \texttt{6}   & \texttt{7}  \\ \hline
            nepřiměřená rychlost jízdy & \texttt{11\,482} & \texttt{1\,424} & \texttt{169} & \texttt{76} \\ \hline
            nesprávné předjíždění      & \texttt{1\,153}  & \texttt{49}   & \texttt{12}  & \texttt{2}  \\ \hline
            nedání přednosti v~jízdě   & \texttt{9\,125}  & \texttt{295}  & \texttt{49}  & \texttt{11} \\ \hline
            nesprávný způsob jízdy     & \texttt{31\,110} & \texttt{3\,369} & \texttt{283} & \texttt{99} \\ \hline
            technická závada vozidla   & \texttt{110}   & \texttt{0}    & \texttt{2}   & \texttt{2}  \\ \hline
        \end{tabular}
    }
\end{table}

Z~analýzy všech záznamů dopravních nehod v~letech 2023 až 2024 (\texttt{163\,580}), které obsahují informace o~hlavních příčinách nehod, průměrné hmotné škodě a~vlivu alkoholu a~drog na viníky nehod, vyplývá, že průměrná hmotná škoda dosahuje hodnoty \texttt{8\,328,15\,Kč}, přičemž jako nejvíce nákladné se ukazují technické závady vozidel s~průměrnou hmotnou škodou \texttt{15\,060,74\,Kč}.
Naopak „nejlevnější“ jsou podle dostupných dat nehody nezaviněné řidičem s~průměrnou hmotnou škodou \texttt{4\,235,26\,Kč}.

Ačkoliv jsou nehody způsobené technickými závadami vozidel nejdražší, jejich celkový počet je ve srovnání s~ostatními příčinami relativně nízký.
Nejčastější příčinou dopravních nehod je však nesprávný způsob jízdy, který má za následek více než pětinu všech kolizí.
Téměř třetina všech dopravních nehod (\texttt{53\,787}) je zaviněna osobami, které nejsou pod vlivem alkoholu či drog.
Navzdory obecným předsudkům ve spojitosti s~agresivitou a~rizikovým chováním účastníků silničního provozu pod vlivem alkoholu a~drog, tito viníci způsobují daleko méně nehod (\texttt{5\,958}) a~jejich průměrná hmotná škoda je překvapivě nižší (\texttt{8\,165,85\,Kč}) než u~střízlivých osob (\texttt{11\,317,65\,Kč}).

\end{document}
